\section{Reduced densities and masses of repeated pure states}
\label{sec:pure_tensor_power}

Let $\mathcal H_S=\mathbb C^S$ and $\mathcal H_R=\mathbb C^\Omega$ be
finite-dimensional, and let $\psi\in\mathcal H_S\otimes\mathcal H_R$.
Write $\rho=\operatorname{tr}_S|\psi\rangle\langle\psi|$ for its actual
reduced density, without imposing normalization.
For $k\ge0$, regroup the repeated vector into the two systems. Its coordinates,
indexed by $x\in S^k$ and $y\in\Omega^k$, are
\begin{align}
 v_k(x,y)=\prod_{i=0}^{k-1}\psi(x_i,y_i).
\end{align}
These are the pure-state identities used in the Schur-sector argument of
\cite[Section~7]{OpenAI2026AreaLaw}.

\begin{theorem}[Reduced density of a repeated bipartite vector]
\label{thm:pure_tensor_power_marginal}
\lean{Matrix.partialTraceLeft_vecMulVec_prod}
\leanok
For every $\psi$ and every $k\ge0$,
\begin{align}
 \operatorname{tr}_{S^k}|v_k\rangle\langle v_k|=\rho^{\otimes k}.
\end{align}
\end{theorem}
\begin{proof}
\leanok
For $y,z\in\Omega^k$, the matrix entry on the left is
\begin{align}
 \sum_{x\in S^k}\prod_i\psi(x_i,y_i)\overline{\psi(x_i,z_i)}
 =\prod_i\sum_{s\in S}\psi(s,y_i)\overline{\psi(s,z_i)}
 =\prod_i\rho(y_i,z_i).
\end{align}
The middle equality is the distributive law for a finite product of sums.
\end{proof}

\begin{theorem}[Projection mass of a repeated bipartite vector]
\label{thm:pure_tensor_power_projection_mass}
\lean{Matrix.norm_sq_one_kronecker_mulVec_prod}
\leanok
\uses{thm:pure_tensor_power_marginal}
For every orthogonal projector $P$ on $\mathcal H_R^{\otimes k}$,
\begin{align}
 \|(I_{S^k}\otimes P)v_k\|^2
 =\operatorname{Re}\operatorname{tr}(\rho^{\otimes k}P).
\end{align}
In particular, taking $P$ to be an actual central Schur projector gives
the squared norm in that sector.
\end{theorem}
\begin{proof}
\leanok
\uses{thm:pure_tensor_power_marginal}
The trace-pairing property of the partial trace and the preceding identity
give
\begin{align}
 \operatorname{tr}(\rho^{\otimes k}P)
 =\langle v_k,(I_{S^k}\otimes P)v_k\rangle
 =\langle (I_{S^k}\otimes P)v_k,(I_{S^k}\otimes P)v_k\rangle.
\end{align}
The second equality uses $P=P^\dagger=P^2$. Taking real parts proves the
formula.
\end{proof}

Both identities include zero copies and empty one-copy factors; no
invertibility of $\rho$ is required. Concentration, the entropy window and
selection of a common sequence of labels are separate assertions.
